% ======================================================================
% Chapter: Logarithmic W-algebras: the triplet family W(p)
% ======================================================================

\chapter{\texorpdfstring{Logarithmic $\cW$-algebras: the triplet
$\cW(p)$}{Logarithmic W-algebras: the triplet W(p)}}
\label{chap:logarithmic-w-algebras}

The non-logarithmic census is controlled by PBW or generic-level
arguments. The triplet algebra $\cW(p)$ is the first finite
non-semisimple test object. Its module category has indecomposable
objects on which $L_0$ has Jordan blocks, yet Zhu's
$C_2$-cofiniteness condition holds, the algebra has finitely many
simple modules, and modular invariance survives in Miyamoto's
pseudo-trace form~\cite{Miyamoto04}. The problem is not whether
$C_2$-cofiniteness restores rationality. It does not. The problem is
which pieces of the modular Koszul package survive for a finite
logarithmic vertex algebra whose $C_2$-algebra is finite-dimensional
and therefore carries normally ordered relations.

The triplet algebra $\cW(p)$ at parameter $p \geq 2$ is the
prototypical $C_2$-cofinite non-rational vertex algebra. It has four
strong generators: the stress tensor~$T$ of weight~$2$ and an
$\mathfrak{sl}_2$-triplet $W^+, W^0, W^-$ of weight $2p - 1$. The
central charge is
\begin{equation}\label{eq:wp-central-charge}
 c(\cW(p)) \;=\; 1 - \frac{6(p-1)^2}{p}.
\end{equation}
At $p = 2$: $c = -2$, the symplectic fermion point. At $p = 3$:
$c = -7$, the first genuinely new case. As $p \to \infty$:
$c \sim -6p$, and the algebra becomes increasingly non-rational.

\begin{table}[ht]
\centering
\small
\caption{Five-theorem interface for
$\cW(p)$.}\label{tab:wp-five-theorems}
\begin{tabular}{lll}
\toprule
\textbf{Theorem} & \textbf{Statement for $\cW(p)$}
 & \textbf{Anchor} \\
\midrule
A (adjunction) &
 Verdier/Koszul dual $\cW(p)^!$ not yet constructed
 & Conj~\ref{conj:wp-koszul-dual} \\
B (inversion) &
 completed bar-cobar counit
 $\Omega(\barB(\cW(p)))\to\cW(p)$
 & Conj~\ref{conj:wp-inversion} \\
C (complementarity) &
 Virasoro-sector sum $13$; full triplet package open
 & Rem~\ref{rem:wp-complementarity} \\
D (modular characteristic) &
 $T$-line scalar $\kappa_T(\cW(p)) = c/2$
 & Prop~\ref{prop:wp-kappa} \\
H (chiral Hochschild) &
 curve-level amplitude $[0,2]$ only under the Koszul-completed
 hypotheses; off-locus degree-$3$ obstruction
 & Prop~\ref{prop:wp-theorem-h-boundary} \\
\bottomrule
\end{tabular}
\end{table}

\begin{table}[ht]
\centering
\small
\caption{Shadow archetype data for the triplet
$\cW(p)$.}\label{tab:wp-shadow-archetype}
\begin{tabular}{ll}
\toprule
\textbf{Invariant} & \textbf{Value} \\
\midrule
Class & M (mixed/infinite shadow depth) \\
Shadow depth $r_{\max}$ & $\infty$ \\
$T$-line $\kappa_T(\cW(p))$ &
 $\dfrac{c(\cW(p))}{2} = \dfrac{1}{2} - \dfrac{3(p-1)^2}{p}$ \\[4pt]
$T$-line $F_1^T$ & $\kappa_T / 24$ \\
Generators & $4$: \; $T$ (wt $2$) $+$ $W^+, W^0, W^-$ (wt $2p{-}1$) \\
$p_{\max}$ (OPE pole order) & $2(2p{-}1) = 4p - 2$ \\
$k_{\max}$ (collision depth) & $4p - 3$ \\
Strong generation & Strong, not free by $C_2$-cofiniteness \\
Koszulness & Open beyond the PBW/free route \\
Complementarity & Virasoro sector known; full triplet package open \\
\bottomrule
\end{tabular}
\end{table}

% ======================================================================
% Section 1: Definition and OPE
% ======================================================================

\section{Definition and OPE structure}\label{sec:wp-definition}

\begin{definition}[Triplet algebra;
\ClaimStatusProvedElsewhere]\label{def:wp-algebra}
\index{triplet algebra|textbf}
\index{logarithmic VOA!triplet}
Let $p \geq 2$ be an integer. The \emph{triplet algebra}
$\cW(p)$ is the vertex subalgebra of the lattice vertex
algebra $V_{\sqrt{2p}\Z}$ generated by the screening kernel:
$\cW(p) = \ker_{\mathrm{res}}(e_0) \cap V_{\sqrt{2p}\Z}$,
where $e_0$ is the zero mode of the screening operator
$Q = \oint e^{\alpha_{+}\phi}$.
Equivalently, $\cW(2) \cong
\mathcal{SF}^{\Z/2\Z}$ (symplectic fermion orbifold),
and more generally $\cW(p)$ is the kernel of the short
screening operator in the rank-$1$ lattice VOA at the
$(2,2p{-}1)$ central charge (Feigin--Ga\u{\i}otto--Frenkel
\cite{FGF2010}, Adamovi\'c--Milas~\cite{AdamovicMilas2008}).

The algebra has \emph{four} strong generators:
\begin{center}
\renewcommand{\arraystretch}{1.3}
\begin{tabular}{cccc}
\toprule
Generator & Conformal weight & $\mathfrak{sl}_2$-charge
 & Parity \\
\midrule
$T$ & $2$ & $0$ & even \\
$W^+$ & $2p-1$ & $+2$ & even \\
$W^0$ & $2p-1$ & $0$ & even \\
$W^-$ & $2p-1$ & $-2$ & even \\
\bottomrule
\end{tabular}
\end{center}
The central charge is $c(\cW(p)) = 1 - 6(p-1)^2/p$.
\end{definition}

\begin{remark}[Triplet and singlet boundary]
\label{rem:wp-singlet-boundary}
\index{singlet algebra!boundary from triplet}
The notation in this chapter is triplet notation. For $p\ge2$,
$\cW(p)$ denotes the $C_2$-cofinite triplet algebra of
Adamovi\'c--Milas: central charge $1-6(p-1)^2/p$, four strong
generators, $2p$ simple ordinary modules, and logarithmic
indecomposables. The singlet algebra $\mathcal M(p)$ has the same
central charge and the same screening origin, but it is not part of
this finite triplet package: the standard singlet literature works
with $C_1$-cofinite finite-length subcategories and with typical
simple modules in continuous families
\textup{(}Creutzig--Ridout~\cite{CreutzigRidout2013}\textup{)}.
No $C_2$-cofinite, finite-simple-module, bar-cobar, or Theorem~H
claim about $\mathcal M(p)$ is imported from the triplet statement.
\end{remark}

\begin{remark}[Generator count and relations]
\label{rem:wp-generator-count}
The $\mathfrak{sl}_2$-triplet $W^+, W^0, W^-$ contributes
\emph{three} strong generators, not one. The total is four,
not two. Since $\cW(p)$ is $C_2$-cofinite, the images of these
four generators in $R_{\cW(p)}=\cW(p)/C_2(\cW(p))$ satisfy
nontrivial polynomial relations. The triplet is finitely strongly
generated but cannot be treated as a freely generated PBW algebra.
\end{remark}

The OPE structure of $\cW(p)$ is determined by three families
of singular terms. The $T$-$T$ OPE is the standard Virasoro OPE
at central charge~\eqref{eq:wp-central-charge}:
\begin{equation}\label{eq:wp-TT-ope}
 T(z)\,T(w) \;\sim\;
 \frac{c/2}{(z-w)^4} + \frac{2T(w)}{(z-w)^2}
 + \frac{\partial T(w)}{z-w}.
\end{equation}
The $T$-$W^a$ OPE records the primary property of the triplet:
\begin{equation}\label{eq:wp-TW-ope}
 T(z)\,W^a(w) \;\sim\;
 \frac{(2p{-}1)\,W^a(w)}{(z-w)^2}
 + \frac{\partial W^a(w)}{z-w},
 \qquad a \in \{+, 0, -\}.
\end{equation}
The $W^a$-$W^b$ OPE has maximal pole order $2(2p{-}1)$:
for the self-OPE $W^0(z)\,W^0(w)$, the leading singularity
is $(z-w)^{-2(2p-1)}$ with coefficient determined by the
central charge and the normalization of $W^0$. The full
$W$-$W$ OPE involves normally ordered composites of $T$
and its descendants through weight $2(2p{-}1) - 2$, together
with the $\mathfrak{sl}_2$-structure constants.

\begin{remark}[OPE pole order and collision depth]
\label{rem:wp-pole-depth}
\index{triplet algebra!pole structure}
The three depth invariants for $\cW(p)$ are:
\begin{align}
 p_{\max}(\cW(p)) &= 2(2p-1) = 4p - 2
 \qquad\text{(max OPE pole order, from $W$-$W$)}, \\
 k_{\max}(\cW(p)) &= 4p - 3
 \qquad\text{(collision depth: $k_{\max} = p_{\max} - 1$)}, \\
 r_{\max}(\cW(p)) &= \infty
 \qquad\text{(shadow depth: class M, from the Virasoro $T$-channel)}.
\end{align}
The $T$-channel alone forces $r_{\max} = \infty$ at the triplet
central charges: the Virasoro shadow tower is non-terminating away
from $c=0$ and $c=-22/5$, and neither value occurs for integer
$p\ge2$. The $W$-channels introduce additional higher-degree shadow
data but do not make the class finite.
\end{remark}

% ======================================================================
% Section 2: The Virasoro-line modular characteristic
% ======================================================================

\section{The Virasoro-line modular characteristic}\label{sec:wp-kappa}

\begin{proposition}[Virasoro-line $\kappa$ for $\cW(p)$;
\ClaimStatusProvedHere]\label{prop:wp-kappa}
\index{triplet algebra!modular characteristic}
The Virasoro $T$-line contribution to the scalar modular
characteristic of the triplet algebra is
\begin{equation}\label{eq:wp-kappa}
 \kappa_T(\cW(p)) \;=\; \frac{c(\cW(p))}{2}
 \;=\; \frac{1}{2} - \frac{3(p-1)^2}{p}.
\end{equation}
The full modular characteristic package of $\cW(p)$ includes the
triplet $W$-channels and is not determined here by the scalar
Virasoro line alone.
\end{proposition}

\begin{proof}
Three checks identify the $T$-line and fence off the full package.

\emph{Path~1 (Virasoro OPE).}
The algebra $\cW(p)$ has $\dim V_1 = 0$: there are no weight-$1$
currents ($T$ has weight~$2$; the triplet $W^a$ has weight
$2p - 1 \geq 3$). The stress tensor nevertheless generates an
autonomous Virasoro subalgebra with
$T(z)T(w)\sim (c/2)(z-w)^{-4}+\cdots$. The ordered collision
residue on this line is therefore the Virasoro residue, and the
averaged scalar is $\kappa_T=c/2$.

\emph{Path~2 (lattice parent).}
The triplet $\cW(p) \subset V_{\sqrt{2p}\Z}$ is a vertex
subalgebra of the rank-$1$ lattice VOA at level $2p$. The
lattice VOA has rank-one Heisenberg scalar $\kappa=1$, but
$\cW(p)$ is a proper screening kernel, not the full lattice VOA.
The Heisenberg current of the parent is not a generator of
$\cW(p)$. The Virasoro sector of $\cW(p)$ has central
charge $c = 1 - 6(p-1)^2/p$, which differs from the lattice value
$c = 1$, and its line contribution is
$c/2=1/2-3(p-1)^2/p$.

\emph{Path~3 (consistency at $p = 2$).}
At $p = 2$: $c(\cW(2)) = 1 - 6/2 = -2$ and
$\kappa_T(\cW(2)) = -1$.
The symplectic fermion algebra $\mathcal{SF}$ (the $bc$ system
at $\lambda = 1$) has $c = -2$
(Computation~\ref{comp:beta-gamma-central-charges}).
The identification
$\cW(2) \cong \mathcal{SF}^{\Z_2}$ (Kausch
\cite{Kausch2000})
preserves the central charge: $c(\cW(2)) = c(\mathcal{SF}) = -2$.
The orbifold preserves the Virasoro vector, hence the $T$-line
scalar is again $\kappa_T=c/2=-1$. It does not follow that the
full triplet modular package is inherited from the parent; the
orbifold introduces non-free invariant relations.
\end{proof}

\begin{remark}[Explicit values]\label{rem:wp-kappa-values}
\begin{center}
\renewcommand{\arraystretch}{1.3}
\begin{tabular}{cccc}
\toprule
$p$ & $c(\cW(p))$ & $\kappa_T(\cW(p))$ & $F_1^T = \kappa_T/24$ \\
\midrule
$2$ & $-2$ & $-1$ & $-1/24$ \\
$3$ & $-7$ & $-7/2$ & $-7/48$ \\
$4$ & $-25/2$ & $-25/4$ & $-25/96$ \\
$5$ & $-91/5$ & $-91/10$ & $-91/240$ \\
$10$ & $-238/5$ & $-119/5$ & $-119/120$ \\
\bottomrule
\end{tabular}
\end{center}
For all $p \geq 2$, $\kappa_T(\cW(p)) < 0$: the triplet
Virasoro line has negative scalar curvature, reflecting the
negative central charge.
\end{remark}


% ======================================================================
% Section 3: C_2-cofiniteness and the bar complex
% ======================================================================

\section{$C_2$-cofiniteness and the bar complex}
\label{sec:wp-c2-cofinite}

The $C_2$-cofiniteness condition
$\dim(\cA / C_2(\cA)) < \infty$
(Zhu~\cite{Zhu96}) is the vertex algebra analogue of
finite-dimensionality for the associated variety
$X_\cA = \operatorname{Spec}(R_\cA)$, where
$R_\cA = \cA / C_2(\cA)$ is Zhu's $C_2$-algebra.
For $C_2$-cofinite algebras, $X_\cA = \{0\}$: the associated
variety is a point.

\begin{proposition}[$C_2$-cofiniteness of $\cW(p)$;
\ClaimStatusProvedElsewhere]\label{prop:wp-c2-cofinite}
\index{triplet algebra!$C_2$-cofiniteness}
For all $p \geq 2$, the triplet algebra $\cW(p)$ is
$C_2$-cofinite. The associated variety
$X_{\cW(p)} = \{0\}$.
\end{proposition}

\begin{proof}[Reference]
Adamovi\'c--Milas~\cite{AdamovicMilas2008}, Theorem~5.1.
The proof uses the explicit screening realization within the
lattice VOA and the $\mathfrak{sl}_2$-action on the
generators.
\end{proof}

\begin{remark}[$C_2$-cofiniteness vs rationality vs Koszulness]
\label{rem:wp-c2-vs-koszul}
\index{$C_2$-cofiniteness!relation to Koszulness}
Three finiteness conditions on a vertex algebra and their
logical relations:
\begin{enumerate}
\item \emph{Rationality} (every admissible module is completely
reducible) $\Longrightarrow$ \emph{$C_2$-cofiniteness} (ABD
theorem, Arakawa--van Ekeren \cite{AvE23}).
\item \emph{$C_2$-cofiniteness} $\not\Longrightarrow$
\emph{rationality}: $\cW(p)$ is $C_2$-cofinite but not
rational for any $p \geq 2$.
\item \emph{Chiral Koszulness}
(Definition~\ref{def:chiral-koszul-morphism}) does not require
free strong generation. Free strong generation is the hypothesis
under which Proposition~\ref{prop:pbw-universality} gives the
fast PBW proof of Koszulness. The triplet $\cW(p)$ is strongly
generated, but it is not freely strongly generated by any finite
positive-weight strong generating set.
\end{enumerate}
The symplectic fermion parent $\mathcal{SF}$ is freely
strongly generated and hence chirally Koszul. Whether the
orbifold $\cW(2) = \mathcal{SF}^{\Z_2}$ inherits Koszulness
from the parent is a non-trivial question: orbifold kernels
do not preserve free strong generation.
\end{remark}

\begin{proposition}[No finite free strong generation;
\ClaimStatusProvedHere]\label{prop:wp-not-free-strong}
\index{triplet algebra!free strong generation}
For $p\ge2$, the triplet algebra $\cW(p)$ is not freely strongly
generated by a finite set of positive conformal-weight fields.
\end{proposition}

\begin{proof}
Let $a_1,\ldots,a_m$ be finite homogeneous strong generators of
positive weights. If they freely strongly generated the vertex
algebra, then the $C_2$-algebra
$R_{\cW(p)}=\cW(p)/C_2(\cW(p))$ would contain the polynomial
monomials in the images $\bar a_1,\ldots,\bar a_m$ as a linearly
independent commutative PBW basis; derivatives disappear in the
$C_2$ quotient, but zero-derivative normally ordered monomials
remain. Thus $R_{\cW(p)}$ would be infinite-dimensional unless all
generators had zero image. The stress tensor and triplet fields are
not all zero in the associated graded of the screening kernel.
Adamovi\'c--Milas prove $C_2$-cofiniteness
(Proposition~\ref{prop:wp-c2-cofinite}), so
$\dim R_{\cW(p)}<\infty$. Hence nontrivial normally ordered
relations must occur, and finite free strong generation is
impossible.
\end{proof}


% ======================================================================
% Section 4: Koszulness: the open problem
% ======================================================================

\section{Koszulness: the central open problem}
\label{sec:wp-koszulness}

\begin{conjecture}[Koszulness of $\cW(p)$;
\ClaimStatusConjectured]\label{conj:wp-koszulness}
\index{triplet algebra!Koszulness conjecture}
The triplet algebra $\cW(p)$ is chirally Koszul for all
$p \geq 2$.
\end{conjecture}

The difficulty is structural, not computational. The PBW/free route
is closed by Proposition~\ref{prop:wp-not-free-strong}; the remaining
routes are quotient, equivariant, and completed routes.

\begin{enumerate}
\item \emph{Screening construction.} The realization
$\cW(p) = \ker(Q) \cap V_{\sqrt{2p}\Z}$ embeds $\cW(p)$
in a freely strongly generated (hence Koszul) lattice VOA.
The kernel of a screening operator on a Koszul algebra
need not be Koszul: the screening kernel can
introduce relations among normally ordered monomials that
destroy the PBW basis.

\item \emph{Presentation by generators and relations.}
Feigin--Tipunin~\cite{FeiginTipunin2010} give an explicit
presentation of $\cW(p)$ by generators and relations.
The presentation establishes finite strong generation and gives the
relations that must be seen by the bar differential. It cannot prove
free strong generation, because $C_2$-cofiniteness already excludes
finite free strong generation.

\item \emph{Character and Zhu comparison.}
The excluded PBW character
$\chi_{\cW(p)}(q) = \prod_{n \geq 2}(1-q^n)^{-1}
\cdot \prod_{n \geq 2p-1}(1-q^n)^{-3}$
is still useful as a defect detector: comparison with the
Adamovi\'c--Milas vacuum character locates the first normally
ordered relation and hence the first candidate bar $2$-cocycle.
It does not by itself prove chiral Koszulness.

\item \emph{Finite-window and completion argument.}
The $T$-line is class~M, so the raw direct-sum class-M
chain-level statement is not the correct ambient. Any positive
Koszulness theorem for $\cW(p)$ must be formulated in finite
conformal-weight windows and then passed to the pro-object /
filtered-completed ambient by a Mittag--Leffler comparison, as in
the repaired class-M higher-genus theorem.
\end{enumerate}

\begin{remark}[Koszulness constraints]
\label{rem:wp-koszul-constraints}
The conjecture is constrained by three finite structures:
\begin{enumerate}
\item The associated variety $X_{\cW(p)} = \{0\}$ is a point. This
is a $C_2$-cofinite constraint on the Li associated graded algebra,
not a Koszulness theorem.
\item The $p = 2$ case $\cW(2) = \mathcal{SF}^{\Z_2}$ has
a Koszul parent, and the equivariant bar complex
(Remark~\ref{rem:symplectic-logarithmic}) detects
the non-semisimplicity at the module level, not the algebra
level.
\item Adamovi\'c--Milas~\cite{AdamovicMilas2008} prove that
$\cW(p)$ has exactly $2p$ simple modules with finite-dimensional
$\Ext$ groups between them. This constrains the finite block
structure of a possible Koszul dual category; it does not determine
the number of algebra generators of~$\cW(p)^!$.
\end{enumerate}
\end{remark}

\begin{proposition}[Theorem~H boundary for the triplet;
\ClaimStatusConditional]\label{prop:wp-theorem-h-boundary}
\index{Theorem H!triplet boundary}
\index{triplet algebra!chiral Hochschild boundary}
Let $p\ge2$, and put $\cA=\cW(p)$ on a smooth projective
reference curve. Assume the conformal-weight completed
Fulton--MacPherson filtration of $RHH_{\mathrm{ch}}(\cA)$ is
exhaustive and Hausdorff, and set
\[
 \mathrm{KD}_{\mathrm H}^{\bullet}(\cA)
 :=F^1_{\mathrm{FM}}RHH_{\mathrm{ch}}(\cA),\qquad
 Q_{\mathrm H}(\cA)
 :=RHH_{\mathrm{ch}}(\cA)/F^1_{\mathrm{FM}}RHH_{\mathrm{ch}}(\cA).
\]
If $Q_{\mathrm H}(\cA)$ has curve-level amplitude $[0,2]$, then
the exact triangle
\[
 \mathrm{KD}_{\mathrm H}^{\bullet}(\cA)
 \longrightarrow RHH_{\mathrm{ch}}(\cA)
 \longrightarrow Q_{\mathrm H}(\cA)
 \xrightarrow{+1}
\]
gives the degree-$3$ obstruction
\[
 \mathfrak{o}_{\mathrm H}^{3}(\cA)
 :=
 \operatorname{coker}\!\left(
  \partial_{\mathrm H}^{3}(\cA)\colon
  H^2(Q_{\mathrm H}(\cA))
  \longrightarrow
  H^3(\mathrm{KD}_{\mathrm H}^{\bullet}(\cA))
 \right)
 \cong \ChirHoch^3(\cA).
\]
For $n\ge4$,
\[
 \ChirHoch^n(\cA)
 \cong H^n(\mathrm{KD}_{\mathrm H}^{\bullet}(\cA)).
\]
Thus a nonzero $\ChirHoch^3(\cW(p))$ is not a fourth column of
curve-level Theorem~H. It is the off-Koszul positive-collision
obstruction $\mathfrak{o}_{\mathrm H}^{3}(\cW(p))$. Under the full
PBW chiral Koszul, completion, perfectness, and averaging hypotheses
of Theorem~\ref{thm:theorem-h-on-koszul-locus}, this obstruction
vanishes and $\ChirHoch^n(\cW(p))=0$ for $n\ge3$.
\end{proposition}

\begin{proof}
This is the specialization of the off-Koszul high-degree exact
sequence of
Theorem~\ref{thm:theorem-h-off-koszul-explicit-correction} to the
triplet algebra. The positive-depth complex is the closed
collision-stratum part of the completed FM Hochschild complex. The
depth-zero quotient supplies only the curve-level $[0,2]$ piece; the
long exact sequence of the displayed triangle identifies the first
possible high-degree class with the cokernel of
$\partial_{\mathrm H}^{3}$. For $\cW(p)$ the first candidates lie on
the $W$-$W$ collision strata of pole order $2(2p-1)$, equivalently
collision depth $4p-3$. Deciding whether these classes die before
$E_\infty$ is exactly the missing Koszul-completed obstruction
calculation.
\end{proof}

% ======================================================================
% Section 5: Non-semisimple module category
% ======================================================================

\section{The non-semisimple module category}
\label{sec:wp-modules}

\begin{proposition}[Module category of $\cW(p)$;
\ClaimStatusProvedElsewhere]\label{prop:wp-modules}
\index{triplet algebra!module category}
The module category $\operatorname{Rep}(\cW(p))$ has exactly
$2p$ simple objects $\Lambda(i)$, $\Pi(i)$ for
$i = 1, \ldots, p$ (Adamovi\'c--Milas~\cite{AdamovicMilas2008}).
It is not semisimple: the non-projective simples have
indecomposable projective covers with Loewy extensions, and the
logarithmic blocks contain modules on which $L_0$ has Jordan
blocks. The category is:
\begin{enumerate}
\item abelian with enough projectives,
\item finite in the logarithmic sense supplied by
$C_2$-cofiniteness (finitely many simples and finite-length
blocks),
\item non-semisimple ($L_0$ has Jordan blocks on indecomposable
projective modules),
\item braided tensor category with non-degenerate Verlinde
formula (Creutzig--Gainutdinov--Runkel~\cite{CreutzigGainutdinovRunkel20}).
\end{enumerate}
\end{proposition}

\begin{remark}[Bar complex and non-semisimplicity]
\label{rem:wp-bar-nonsemisimple}
\index{bar complex!non-semisimple modules}
The non-semisimplicity of $\operatorname{Rep}(\cW(p))$ is
visible in the bar complex through two mechanisms:
\begin{enumerate}
\item \emph{Jordan blocks generate non-trivial $d_2$
differentials} in the bar spectral sequence
(Remark~\ref{rem:jordan-bar-spectral}): a length-$r$
Jordan block for $L_0$ acting on an indecomposable module
$P(i)$ produces non-trivial $E_r$ differentials connecting
composition factors.
\item \emph{The Ext algebra is non-semisimple}: the
finite-dimensional $\Ext^*(\Lambda(i), \Lambda(j))$ groups
(Proposition~\ref{prop:wp-modules}) are detected by bar
$1$-cocycles, and their non-vanishing higher products encode
the non-trivial $A_\infty$-structure on the derived category.
\end{enumerate}
If $\cW(p)$ is chirally Koszul, these structures are organized
by the Koszul dual $\cW(p)^!$, whose representation category
controls the $\Ext$-algebra. The non-semisimplicity would then
be a \emph{module-level} phenomenon, compatible with
\emph{algebra-level} Koszulness, exactly as for the $\beta\gamma$
parent (Remark~\ref{rem:symplectic-logarithmic}).
\end{remark}


% ======================================================================
% Section 6: Shadow depth and the G/L/C/M classification
% ======================================================================

\section{Shadow depth and the G/L/C/M classification}
\label{sec:wp-shadow-depth}

\begin{proposition}[Shadow class of $\cW(p)$;
\ClaimStatusConditional]\label{prop:wp-shadow-class}
\index{triplet algebra!shadow class}
The triplet algebra $\cW(p)$ is class~M for all $p \geq 2$:
$r_{\max}(\cW(p)) = \infty$.
\end{proposition}

\begin{proof}
The shadow obstruction tower is governed channel-by-channel.
The $T$-channel of $\cW(p)$ is a copy of the Virasoro OPE
at central charge $c = 1 - 6(p-1)^2/p$. By the single-line
dichotomy (Theorem~\ref{thm:single-line-dichotomy}), the
$T$-channel has critical discriminant
$\Delta_T = 8\kappa_T S_4^T = 40/(5c+22)$. For integer
$p\ge2$, neither $c=0$ nor $c=-22/5$ occurs. Therefore the
$T$-channel shadow metric $Q_T(t)$ is irreducible, and the
$T$-channel tower does not terminate: $r_{\max}^T = \infty$.
Since the triplet tower contains this non-terminating regular
$T$-line subtower, the overall shadow depth is infinite.
\end{proof}

\begin{remark}[Multi-channel structure]
\label{rem:wp-multichannel}
The shadow obstruction tower of $\cW(p)$ has a
multi-channel structure similar to the $\mathcal{N}=2$
superconformal algebra
(Table~\ref{tab:n2-shadow-archetype}):
\begin{center}
\renewcommand{\arraystretch}{1.3}
\begin{tabular}{lcl}
\toprule
Channel & Class & Depth \\
\midrule
$T$-$T$ & M & $\infty$ \\
$T$-$W^a$ & & $\geq 3$ (primary OPE) \\
$W^a$-$W^b$ & & $\geq 4p - 3$ \\
\bottomrule
\end{tabular}
\end{center}
The $W$-channels have high collision depth $k_{\max} = 4p - 3$
(from the $(4p{-}2)$-fold pole in the $W$-$W$ OPE, with
$d\log$ absorption reducing by one), which generates shadow
data at correspondingly high degrees. The overall class is
determined by the $T$-channel, which forces $r_{\max} = \infty$.
\end{remark}


% ======================================================================
% Section 7: The Koszul dual: conjectures
% ======================================================================

\section{The Koszul dual and complementarity}
\label{sec:wp-koszul-dual}

\begin{conjecture}[Koszul dual of $\cW(p)$;
\ClaimStatusConjectured]\label{conj:wp-koszul-dual}
\index{triplet algebra!Koszul dual}
If the triplet algebra is chirally Koszul, its Verdier/Koszul
dual $\cW(p)^!$ exists in the completed chiral category. Its
Virasoro sector has central charge
$c' = 26 - c(\cW(p)) = 25 + 6(p-1)^2/p$ and
$T$-line scalar
$\kappa_T' = c'/2 = 25/2 + 3(p-1)^2/p$.
No assertion is made here that the full dual is again a triplet
algebra, a $C_2$-cofinite algebra, or a lattice screening kernel.
\end{conjecture}

\begin{remark}[Five-object firewall for the triplet]
\label{rem:wp-five-object-firewall}
The symbols in the previous conjecture denote different objects:
$\cW(p)$ is the algebra; $\barB(\cW(p))$ is its bar coalgebra;
$H^\bullet(\barB(\cW(p)))$, when it is concentrated in the Koszul
way, is the intrinsic bar-dual coalgebra $\cW(p)^i$ rather than an
algebra; $\cW(p)^!$ is obtained only after the Verdier finite-dual
step $\mathbb D_X(\cW(p)^i)$. The derived chiral centre
$Z^{\mathrm{der}}_{\mathrm{ch}}(\cW(p))$ is the Hochschild cochain
object. The counit
$\Omega(\barB(\cW(p)))\to\cW(p)$ is inversion in the completed
bar-cobar adjunction, not the construction of $\cW(p)^!$.
A physical bulk algebra is identified with
$Z^{\mathrm{der}}_{\mathrm{ch}}(\cW(p))$ only after a named
open/closed comparison quasi-isomorphism; no such comparison is
constructed here.
\end{remark}

\begin{remark}[Complementarity for the Virasoro sector]
\label{rem:wp-complementarity}
For the $T$-channel (Virasoro sector), the complementarity
sum is the standard Virasoro value:
\begin{equation}\label{eq:wp-complementarity}
 \kappa_T(\cW(p)) + \kappa_T(\cW(p)^!)
 \;=\; \frac{c}{2} + \frac{26 - c}{2}
 \;=\; 13.
\end{equation}
This is the Virasoro-sector complementarity. The full
complementarity involving the $W$-channels is open, and
determining it requires knowing whether the Koszul dual
$\cW(p)^!$ has an $\mathfrak{sl}_2$-triplet structure at a
shifted weight.
\end{remark}

\begin{conjecture}[Bar-cobar inversion;
\ClaimStatusConjectured]\label{conj:wp-inversion}
If $\cW(p)$ is chirally Koszul, then the bar-cobar counit
$\Omega(\barB(\cW(p))) \to \cW(p)$
is a quasi-isomorphism after conformal-weight completion, equivalently
stagewise on every finite conformal-weight window with
Mittag--Leffler passage to the pro-object limit.
\end{conjecture}


% ======================================================================
% Section 8: The p=2 case: connection to symplectic fermions
% ======================================================================

\section{The case \texorpdfstring{$p = 2$}{p = 2}:
symplectic fermions}\label{sec:wp-p2}

At $p = 2$ the triplet algebra is $\cW(2) \cong
\mathcal{SF}^{\Z_2}$, the $\Z_2$-orbifold of the
symplectic fermion system at $c = -2$. The strong generators
are $T$ (weight~$2$) and the $\mathfrak{sl}_2$-triplet
$W^+, W^0, W^-$ (weight~$3$).

\begin{computation}[Invariants of $\cW(2)$]
\label{comp:wp-p2-invariants}
\begin{align}
 c(\cW(2)) &= 1 - 6/2 = -2, \\
 \kappa_T(\cW(2)) &= -1, \\
 F_1^T(\cW(2)) &= -1/24, \\
 p_{\max} &= 6 \quad\text{($W$-$W$ OPE)}, \\
 k_{\max} &= 5.
\end{align}
\end{computation}

The symplectic fermion parent $\mathcal{SF}$ has $c = -2$ and
$T$-line scalar $\kappa_T(\mathcal{SF}) = c/2 = -1$, and is
chirally Koszul.
The $\Z_2$-orbifold $\cW(2) = \mathcal{SF}^{\Z_2}$ preserves
the central charge and the Virasoro vector. It does not preserve
free strong generation, and the full triplet modular package is not
the parent package. Whether Koszulness persists through the orbifold
is precisely
the content of Conjecture~\ref{conj:wp-koszulness} at
$p = 2$; the equivariant bar analysis of
Remark~\ref{rem:symplectic-logarithmic} is the relevant
comparison.

\begin{remark}[Module category at $p = 2$]
\label{rem:wp-p2-modules}
$\cW(2)$ has $4 = 2p$ simple modules:
$\Lambda(1)$ ($h = 0$), $\Lambda(2)$ ($h = 1$),
$\Pi(1)$ ($h = -1/8$), $\Pi(2)$ ($h = 3/8$)
in the Adamovi\'c--Milas labelling.
The Grothendieck ring
$K_0(\operatorname{Rep}(\cW(2))) \cong \Z^4$ with a
non-semisimple tensor product reflecting the Jordan block
structure.
\end{remark}


% ======================================================================
% Section 9: Higher p: the general picture
% ======================================================================

\section{Higher \texorpdfstring{$p$}{p}: the general
picture}\label{sec:wp-higher-p}

As $p$ increases, the triplet algebra grows in complexity:

\begin{enumerate}
\item The weight of the triplet generators increases as
$2p - 1$, making the $W$-$W$ OPE pole order $4p - 2$
progressively deeper. The bar complex at degree~$2$
involves modes through order $4p - 3$.

\item The number of simple modules grows as $2p$, and the
non-semisimple structure of the module category becomes
more elaborate: the projective covers $P(i)$ have
socle filtrations of increasing length.

\item The Virasoro-line scalar
$\kappa_T(\cW(p)) = 1/2 - 3(p-1)^2/p
= -3p + 13/2 - 3/p$
grows linearly negative in $p$. The $T$-line genus-$1$ free
energy $F_1^T = \kappa_T/24$ is correspondingly negative,
reflecting the large negative central charge on that line.

\item The Verlinde formula
(Creutzig--Gainutdinov--Runkel~\cite{CreutzigGainutdinovRunkel20}) for
the $\cW(p)$ modular tensor category involves
generalized $S$-matrices compatible with
pseudo-trace functions, not ordinary traces.
The corresponding shadow problem is to locate the pseudo-trace
contribution inside the positive collision-depth filtration.
\end{enumerate}

\begin{remark}[Relation to $(p,1)$ minimal models]
\label{rem:wp-minimal-models}
The $(p,1)$ minimal model has $c = 1 - 6(p-1)^2/p$,
the same central charge as $\cW(p)$. But the $(p,1)$
minimal model is the \emph{simple quotient}
$L(c_{p,1}, 0)$, while $\cW(p)$ is a larger algebra
containing additional generators. The triplet algebra
sits between the Virasoro algebra and the lattice VOA:
$\mathrm{Vir}_c \subset \cW(p) \subset V_{\sqrt{2p}\Z}$.
The $(p,1)$ model is the Virasoro subalgebra's contribution,
while $\cW(p)$ includes the full screening kernel.
\end{remark}


% ======================================================================
% Section 10: Open problems
% ======================================================================

\section{Open problems}\label{sec:wp-open-problems}

\begin{enumerate}
\item \textbf{The first non-free relation.}
Finite free strong generation is impossible by
Proposition~\ref{prop:wp-not-free-strong}. The open problem is to
identify the first normally ordered relation in the
Feigin--Tipunin presentation and compute the corresponding bar
$2$-cocycle. Its vanishing or survival is the first concrete test of
Conjecture~\ref{conj:wp-koszulness}.

\item \textbf{The Koszul dual $\cW(p)^!$.}
Construct the Koszul dual explicitly. Is it again a
$C_2$-cofinite algebra? Does it have a lattice VOA
realization? Is $\operatorname{Rep}(\cW(p)^!)$
non-semisimple?

\item \textbf{Logarithmic Verlinde and the shadow CohFT.}
The generalized Verlinde formula for $\cW(p)$ involves
pseudo-trace functions (Miyamoto~\cite{Miyamoto04}).
Determine whether the shadow CohFT
(Theorem~\ref{thm:shadow-cohft}) extends to the
logarithmic setting, with pseudo-trace functions
replacing ordinary traces.

\item \textbf{Bar complex for non-freely-generated algebras.}
Develop the bar complex theory for vertex algebras that
are strongly but not freely strongly generated. The
conjectural model places the relations among normally ordered
monomials, in the completed or filtered bar construction, as bar
$2$-cocycles rather than as a failure of $C_2$-finiteness.

\item \textbf{Higher-rank logarithmic algebras.}
The triplet $\cW(p)$ is a rank-$1$ logarithmic algebra.
Higher-rank analogues (the $\cW(p,q)$ algebras of
Feigin--Ga\u{\i}otto--Frenkel~\cite{FGF2010}) have
richer structure and more elaborate non-semisimple
module categories. Determine their shadow class and
modular characteristics.

\item \textbf{The triple Massey class
$\langle \Omega, \Omega, \Omega \rangle$: two distinct objects.}
\label{item:wp-massey-omega-triple}
Since $\cW(p)$ fails finite free strong generation, the candidate
first bar-cohomology obstruction is recorded as a triple Massey
product of the augmentation element
$\Omega \in \bar\cW(p) = \ker(\epsilon)$ in
$H^\bullet(B(\cW(p)))$. Two distinct Massey objects live on
$\cW(p)$ and must be distinguished.

The \emph{correlation-function Massey} is the secondary
cohomology operation on the logarithmic $n$-point amplitudes
$\langle V_1(z_1) \cdots V_n(z_n) \rangle$ of $\cW(p)$. The
indecomposable $L_0$-Jordan partners of
Proposition~\ref{prop:wp-modules} produce
$(\log z_{ij})^n$-enriched amplitudes whose Massey-product
norm grows unboundedly in the graph vertex count.
Gurarie~\cite{Gurarie1993} and Flohr~\cite{Flohr1996}
analyse this object and establish its unboundedness despite
$C_2$-cofiniteness.

The \emph{shadow-tower Massey} is the secondary operation on
the structure constants $S_r$ of the bar tower, extracted by
the Arnold $d\!\log$-residue against the Fulton--MacPherson
compactification $\mathrm{FM}_r(\CC)$. It lives in the
polynomial-$d\!\log$ subcomplex
$H^\bullet_{\mathrm{reg}}(B(\cW(p)))$, the homogeneous image
of the Orlik--Solomon grading. The $(\log z_{ij})^n$ with
$n \geq 2$ sit in the kernel of the homogeneous residue by
construction, so the shadow tower does not see the
logarithmic sector at all. The two sectors are orthogonal
under the Orlik--Solomon pairing; the two Massey products
are distinct secondary operations on distinct subcomplexes.

The implication
``$C_2$-cofiniteness $\Rightarrow$ bounded correlation-function
Massey'' is false by Gurarie--Flohr. The implication
``$C_2$-cofiniteness $\Rightarrow$ bounded shadow-tower
Massey'' is open; the sharpest test case is $p = 3$ (first
non-symplectic-fermion), where the shadow-tower coefficients
$S_r(\cW(3))$ are computed through $r = 10$ by the $T$-line
and Cartan-line recursions. The explicit cocycle representative
of $\langle \Omega, \Omega, \Omega \rangle_{\mathrm{shadow}}$
in the bar complex of $\cW(3)$ is not inscribed in this
chapter; its isolation is recorded as
Conjecture~\ref{conj:wp-shadow-tower-massey-p3} below.
\end{enumerate}

\begin{conjecture}[Shadow-tower triple Massey bounded on $\cW(3)$;
\ClaimStatusConjectured]\label{conj:wp-shadow-tower-massey-p3}
\index{triplet algebra!shadow-tower Massey bound}
The triple Massey class
$\langle \Omega, \Omega, \Omega \rangle_{\mathrm{shadow}}
\in H^3_{\mathrm{reg}}(B(\cW(3)))$
of the augmentation element $\Omega$ in the regular sector of
the bar cohomology of the triplet algebra at $p = 3$ admits a
cocycle representative whose Arnold norm is bounded by a
constant depending only on the finite Zhu algebra
$A(\cW(3))$, the six simple-module characters, and the central
charge $c = -7$,
independent of the bar weight stratum. Equivalently, the
shadow coefficients $S_r^T(\cW(3))$ on the $T$-line and
$S_r^{\mathrm{Cartan}}(\cW(3))$ on the Cartan line satisfy
\[
\limsup_{r \to \infty}
|S_r^T(\cW(3))|^{1/r} < \infty,
\qquad
\limsup_{r \to \infty}
|S_r^{\mathrm{Cartan}}(\cW(3))|^{1/r} < \infty.
\]
\end{conjecture}

\begin{remark}[Finite-window check at $p = 3$]
\label{rem:wp-massey-check-p3}
The $T$-line of $\cW(3)$ carries the Virasoro shadow tower at
$c = -7$. The Virasoro shadow asymptotic base
$C_{\Vir_c} = 6$ of
Proposition~\ref{prop:universal-base-CA-six} is
$c$-independent on the non-logarithmic principal stratum;
specialising to $c = -7$ gives a well-defined finite base on
the $T$-line subsector. The Cartan line carries a distinct
quadratic shadow coefficient
$S_2^{\mathrm{Cart}}=1/(2h_W)=1/(2(2p-1))$, which at
$p = 3$ equals $1/10$
(Proposition~\ref{prop:wp-cartan-shadow-through-r6}). The product
of the two asymptotic bases gives a finite upper bound on the
mixed-channel shadow growth through the available window
$r \leq 10$;
extrapolation to all $r$ is the content of the conjecture.
\end{remark}

\begin{remark}[Falsification test]
\label{rem:wp-massey-falsification-p3}
The conjecture falsifies as follows. Compute three further
shadow coefficients $S_{11}, S_{12}, S_{13}$ at $p = 3$ on
both the $T$-line and the Cartan line from the same recurrence. Form
the Stirling-regularised sequence
$|S_r|^{1/r} \cdot (r!)^{-1/r}$ and check against the
Virasoro baseline $C_{\Vir_{-7}} = 6$. Divergence of the
ratio within the extended window falsifies
Conjecture~\ref{conj:wp-shadow-tower-massey-p3}; stabilisation
at a finite constant near the Virasoro base confirms it. Raw
$|S_r|^{1/r}$ is consistent with either hypothesis in a
finite window; the Stirling-regularised ratio is the sharp
test.
\end{remark}

\begin{remark}[Universal $p$ extension is further conjectural]
\label{rem:wp-massey-universal-p}
The analogous statement at all $p \geq 2$
(shadow-tower triple Massey bounded uniformly in $p$) is
not a consequence of $C_2$-cofiniteness of $\cW(p)$
for every $p \geq 2$. The uniformity of the bound in $p$
is not implied by the $p = 3$ case. Adamovi\'c--Milas
\cite{AdamovicMilas2008} supply the simple-module count $2p$
and explicit simple-module characters. The finite Zhu algebra
contains additional nilpotent data and is not measured by the
simple count alone. These character data do not directly produce a
shadow-tower Massey bound: chiral transport from character to
Orlik--Solomon residue requires a separate lemma
identifying the homogeneous residue with a finite-dimensional
pseudo-trace. The $p = 3$ case is the primitive. The
universal-$p$ extension and the full shadow-tower
Koszulness implication are precisely
Conjecture~\ref{conj:wp-koszulness}.
\end{remark}
